% % \begin{table*}[t]
%     % \setlength{\abovecaptionskip}{-0.cm}
%     \setlength{\belowcaptionskip}{-0.2cm}
%     % \renewcommand{\arraystretch}{1.2}
%     \Large
%     \centering
%     \resizebox{0.9\textwidth}{!}{
%     \begin{tabular}{l|ccc|ccc}
%         \toprule
%         ~ & \multicolumn{3}{c}{\textbf{Content}} & \multicolumn{3}{c}{\textbf{Timbre}}  \\
%         ~ & \multicolumn{2}{c}{DSRIBench} & Reconstruction &
%         \multicolumn{2}{c}{DSRIBench} & Reconstruction \\
%         ~ & MI & WER & WER& MI & ACC & SIM \\
%         \midrule
%         EnCodec-L1 & 16.5 & 61.52 & 38.34 & 6.52 & 0.44 & 0.92 \\
%         EnCodec-All & 23.6 & 30.91 & 5.11 & 14.7 & 0.94 & 0.98 \\
%         \method-L1 & & & 9.57 & & & 0.74 \\
%         \method-L2:8 & & & 86.58 & & & 0.79 \\
%         \method-all & & & 5.04 & & & 0.97 \\
%         \bottomrule
%     \end{tabular}}
%     \caption{Results of speech disentanglement.
%     }
%     \label{tab:speech_disentanglement}
% \end{table*}


\begin{table*}[h]
    % \setlength{\abovecaptionskip}{-0.cm}
    \setlength{\belowcaptionskip}{-0.2cm}
    % \renewcommand{\arraystretch}{1.2}
    \Large
    \centering
    \resizebox{0.8\textwidth}{!}{
    \begin{tabular}{lc|cc|cc}
        \toprule
        ~ & ~& \multicolumn{2}{c}{\textbf{TA}} & \multicolumn{2}{c}{\textbf{IP}}   \\
        Teacher &  & MI & WER$^{\dagger}$ & WER$^{\ast}$ & SIM\\
        \midrule
        HuBERT L9 unit & RVQ-1 &  &  & 20.02 & 0.72 \\
        HuBERT L9 unit & RVQ-8 &  &  & 5.84 & 0.95 \\
        HuBERT L9 representation & RVQ-1 & & & & \\
        HuBERT L9 representation & RVQ-8 & & & & \\
        HuBERT average representation & RVQ-1 &  &  & 9.57 & 0.74 \\
        HuBERT average representation & RVQ-8 &  &  & 5.04 & 0.97 \\
        \bottomrule
    \end{tabular}}
    \caption{Results of \method training with different semantic teacher representations on DSRBench. TA refers to text alignment and IP denotes information preservation. MI and WER$^{\dagger}$ refer to mutual information and word error rate of the downstream model. WER$^{\ast}$ and SIM refer to word error rate and speaker similarity of resynthesized speech respectively.
    }
    \label{tab:teacher_ablation}
\end{table*}

As shown in Table~\ref{tab:bench_results}, as semantic teachers, HuBERT L9 representations perform better than HuBERT units in both Text Alignment and Information Preservation, regardless of whether it's RVQ-1 or RVQ-1:8. The reason may be that discrete HuBERT units lose some content information compared to the continuous representations, thereby providing weaker semantic guidance to \method.

When comparing HuBERT L9 representations with HuBERT average representations, we find that in terms of Text Alignment, the mutual information is higher when HuBERT L9 representations serve as the teacher. This is because HuBERT average representations contain some timbre information, while HuBERT L9 offers purer content information. On the other hand, HuBERT average shows better performance in Information Preservation, reflected in a lower WER. We speculate that this is due to a certain level of task conflict between semantic distillation and reconstruction, where the former aims to retain only content information while the later aims to preserve various aspects of speech. The presence of some timbre information in HuBERT average representations could to some extent alleviate this task conflict.
% 比较HuBERT units和HuBERT L9 representations分别作为semantic teacher时在\bench上的表现，可以发现不管是RVQ-1还是RVQ-1:8，HuBERT L9 representations的在Text Alignment和Information Preservation上都要更好一些。我们推测原因可能是离散表示HuBERT unit相比于连续表示会损失掉一部分内容信息，从而给\method带来的semantic guidance不够强。
% 比较HuBERT L9 representation和HuBERT average representation，我们发现在Text Alingment上，HuBERT L9 representation作为teacher的mutual information要更高。这是因为HuBERT average representation会含有一些timbre information，HuBERT L9的content information纯度更高。而在Information Preservation上HuBERT L9的表现要更好一些，有更低的WER，我们猜测是因为semantic distillation的目标是要只保留内容信息除去其他信息，reconstruction的目标是保留语音中各方面信息，他们之间存在一定的多任务冲突，而HuBERT average representation中含有少量的timbre information，会在一定程度上缓解这种多任务冲突。